\section{Discussion}

\subsection{Interpretation of Results}

Our comprehensive evaluation reveals nuanced performance differences across MT tools and text categories. The results challenge the assumption that a single "best" tool exists for all software localization scenarios.

\subsubsection{Technical Documentation Translation}

[Tool X]'s superior performance on technical documentation (BLEU=X.XX) can be attributed to:
\begin{itemize}
    \item Better handling of complex sentence structures
    \item More consistent technical terminology
    \item Superior context understanding in longer texts
\end{itemize}

However, the XX\% cost premium may not be justified for all projects, particularly when [Tool Y] achieves XX\% of the quality at XX\% of the cost.

\subsubsection{UI String Translation Challenges}

The lower scores on UI strings (average BLEU=X.XX vs X.XX for technical docs) reflect inherent challenges:

\begin{enumerate}
    \item \textbf{Brevity}: Short texts provide limited context for disambiguation
    \item \textbf{Imperatives}: Turkish imperative mood requires careful handling
    \item \textbf{Formality}: UI strings require appropriate formality level
    \item \textbf{Placeholders}: Variable preservation is critical but challenging
\end{enumerate}

\subsection{Practical Implications}

\subsubsection{Tool Selection Recommendations}

Based on our findings, we recommend:

\begin{itemize}
    \item \textbf{For Technical Documentation}: Use [Tool X] when quality is paramount, or [Tool Y] for balanced cost-quality
    \item \textbf{For UI Strings}: [Tool Y] provides best results; consider manual review for critical strings
    \item \textbf{For Error Messages}: [Tool X] ensures accuracy; invest in human review given criticality
    \item \textbf{For High-Volume Projects}: Microsoft Translator offers best cost efficiency
\end{itemize}

\subsubsection{Best Practices}

\begin{enumerate}
    \item \textbf{Provide Context}: When possible, translate entire paragraphs rather than isolated sentences
    \item \textbf{Use Glossaries}: All tools support custom terminology; leverage this for consistency
    \item \textbf{Post-Edit Critical Content}: Allocate human review budget to error messages and key UI elements
    \item \textbf{Test Placeholders}: Always verify placeholder preservation in localized software
    \item \textbf{Consider Hybrid Approaches}: Use different tools for different content types
\end{enumerate}

\subsection{Limitations}

Our study has several limitations:

\begin{enumerate}
    \item \textbf{Language Pair}: Results specific to English-Turkish; other pairs may differ
    \item \textbf{Automatic Metrics}: While correlated with human judgment, automatic metrics don't capture all quality aspects
    \item \textbf{Temporal Validity}: MT systems continuously improve; results reflect [Date] versions
    \item \textbf{Domain Specificity}: Focused on software domain; other technical domains may show different patterns
    \item \textbf{No Fine-Tuning}: We evaluated default models; custom-trained models might perform differently
\end{enumerate}

\subsection{Comparison with Related Work}

Our findings align with previous WMT evaluations \cite{wmt2020findings} showing DeepL's strong performance on European languages. However, our domain-specific focus reveals performance variations not apparent in general-domain evaluations.

The lower scores on UI strings compared to technical documentation (XX\% decrease) corroborate findings from Mozilla's localization studies \cite{mozilla2018localization}, emphasizing the unique challenges of translating short, context-dependent strings.

\subsection{Implications for Research}

Our results suggest several directions for future research:

\begin{enumerate}
    \item \textbf{Context-Aware Translation}: Developing methods to provide UI context to MT systems
    \item \textbf{Placeholder-Aware Models}: Training models specifically to preserve formatting elements
    \item \textbf{Domain Adaptation}: Fine-tuning general MT models on software localization data
    \item \textbf{Turkish-Specific Optimization}: Addressing morphological challenges in Turkish generation
    \item \textbf{Hybrid Systems}: Combining strengths of different tools
\end{enumerate}

\subsection{Reproducibility}

To ensure reproducibility, we provide:
\begin{itemize}
    \item Complete dataset with preprocessing scripts
    \item API wrapper implementations
    \item Evaluation metric calculations
    \item Statistical analysis notebooks
    \item Web application for interactive exploration
\end{itemize}

All materials are available at [GitHub URL] under [License].
